% Generated draft: requires mathematical and source review.
\documentclass[11pt]{article}
\usepackage[margin=25mm]{geometry}
\usepackage{fontspec}
\setmainfont{Latin Modern Roman}
\usepackage{amsmath,amssymb,mathtools,unicode-math}
\setmathfont{Latin Modern Math}
\usepackage{enumitem,needspace,xurl,hyperref}
\setlength{\parindent}{0pt}
\setlength{\parskip}{5pt}
\setlength{\emergencystretch}{3em}
\newcommand{\N}{\mathbb{N}}
\newcommand{\Z}{\mathbb{Z}}
\newcommand{\Q}{\mathbb{Q}}
\newcommand{\R}{\mathbb{R}}
\newcommand{\C}{\mathbb{C}}
\newcommand{\F}{\mathbb{F}}
\newcommand{\A}{\mathbb{A}}
\newcommand{\PP}{\mathbb{P}}
\newcommand{\E}{\mathbb{E}}
\newcommand{\eps}{\varepsilon}
\newcommand{\dd}{\,\mathrm{d}}
\begin{document}
\section*{UnsolvedMath ID: 1863}
\textbf{Source problem code:} catalogue entry
\subsection{Definitions and mathematical statement}
Let \(r\geq1\), and let \(H=\{h_1,\ldots,h_r\}\subset\Z_{\geq0}\) be a finite offset set with \(0=h_1<h_2<\cdots<h_r\). Call \(H\) admissible if for every prime \(\ell\), its residues do not cover all of \(\Z/\ell\Z\):
\[
|\{h_i\bmod\ell:1\leq i\leq r\}|<\ell.
\]
This is the usual exact-offset interpretation of ``no congruence obstructions.'' Say that a positive integer \(t\) realizes \(H\) as consecutive primes if
\[
\{q:q\text{ prime},\ t\leq q\leq t+h_r\}
=\{t+h_1,\ldots,t+h_r\}.
\]
Thus every prescribed entry is prime and there is no additional prime between the first and last entries.

Under these explicit interpretations, the source question asks whether
\[
\forall r\geq1\ \forall\text{ admissible }H=\{0=h_1<\cdots<h_r\},
\quad
\#\{t\in\Z_{>0}:t\text{ realizes }H\text{ as consecutive primes}\}=\infty.
\]
The meanings of pattern, obstruction and consecutiveness are not defined in the source statement itself; the formula is the natural admissible exact-offset reading used by the supplied triage and still requires confirmation against Guy's original formulation. In particular the pattern \(\{0,2\}\) is included, and the background's triples \(\{6k-1,6k+1,6k+5\}\) and \(\{6k+1,6k+5,6k+7\}\) correspond to offsets \(\{0,2,6\}\) and \(\{0,4,6\}\).

The supplied triage dated 2026-08-17 labels this reading partially solved while reporting no universal all-pattern theorem. It distinguishes Maynard's bounded clusters of every fixed size and positive-proportion approximations from realizing every prescribed tuple; Banks--Freiberg--Turnage-Butterbaugh's consecutive primes in suitable tuples and structured increasing, decreasing and divisibility gap patterns; and Pintz's existence, for each size, of some bounded recurring consecutive-prime patterns whose translation sets contain arbitrarily long finite arithmetic progressions. None changes the universal prescribed-pattern quantifier above to an established theorem. The triage also records that, under generalized Elliott--Halberstam, the cited Polymath result gives at least two prime entries in translates of any admissible triple, rather than all three.

For historical context, the background attributes to Hensley--Richards an incompatibility with the assertion
\[
\pi(x+y)\leq\pi(x)+\pi(y)\qquad(x,y\in\Z,\ x,y\geq2),
\]
where \(\pi(u)\) counts primes at most \(u\). The triage verifies the bibliographic identity of their paper but does not use or independently certify that historical incompatibility claim. Its phrase ``more general than Chowla's conjecture'' is historically ambiguous because no exact Chowla formulation is specified. The registry status is open; all progress and status descriptions here are supplied dated assessments.
\subsection{Short English statement}
Does every finite admissible exact pattern of offsets occur infinitely often as an uninterrupted block of consecutive primes? This interpretation includes twin primes. The supplied progress finds some bounded patterns or only some prime entries in a tuple, rather than proving the claim for every prescribed pattern.
\subsection{Sources}
\begin{itemize}
\item Local UnsolvedMath registry, permanent ID 1863, source code catalogue entry, ``General Patterns of Consecutive Primes,'' including supplied literature triage dated 2026-08-17. The background attributes the source to Richard Guy, \emph{Unsolved Problems in Number Theory}, Section A9; no edition or page is supplied.
\item James Maynard, ``Small gaps between primes,'' \emph{Annals of Mathematics} 181 (2015), 383--413, DOI 10.4007/annals.2015.181.1.7: \url{https://annals.math.princeton.edu/2015/181-1/p07}.
\item William D. Banks, Tristan Freiberg, and Caroline L. Turnage-Butterbaugh, ``Consecutive primes in tuples,'' \emph{Acta Arithmetica} 167 (2015), 261--266, DOI 10.4064/aa167-3-4; arXiv:1311.7003: \url{https://arxiv.org/abs/1311.7003}.
\item Janos Pintz, ``Patterns of primes in arithmetic progressions,'' arXiv:1509.01564v2 (2015): \url{https://arxiv.org/abs/1509.01564}.
\item D. H. J. Polymath, ``Variants of the Selberg sieve, and bounded intervals containing many primes,'' \emph{Research in the Mathematical Sciences} 1 (2014), article 12; arXiv:1407.4897: \url{https://arxiv.org/abs/1407.4897}.
\item Douglas Hensley and Ian Richards, ``Primes in intervals,'' \emph{Acta Arithmetica} 25 (1974), 375--391, DOI 10.4064/aa-25-4-375-391: \url{https://www.impan.pl/en/publishing-house/journals-and-series/acta-arithmetica/all/25/4/100146/primes-in-intervals}.
\end{itemize}
\end{document}
